\documentclass[12pt]{article}
\usepackage[T1]{fontenc}
\usepackage[utf8]{inputenc}
\usepackage{lmodern}
\usepackage{textcomp}
\usepackage{enumitem}
\usepackage{geometry}
\geometry{margin=1in}
\begin{document}
\begin{center}
Answer Key - Mathematics - Graduate Level Assessment 15 \
\small (Answers are ordered exactly as in the question paper.)
\end{center}

\section*{Multiple Choice}
\begin{enumerate}[label=\arabic*.]
\item \textbf{Answer:} B
\\ \textit{Options:} A) True, False; B) True, True; C) False, False; D) False, True
\\ \textit{Note:} correct because the formula for the order of the product of two subgroups, |HK| = |H||K|/|H ∩ K|, is a standard result in group theory, and the classification of groups of order 2 p, where p is an odd prime, establishes that they are either cyclic or isomorphic to the dihedral group D\_p, which aligns with the properties of groups in this order.
\item \textbf{Answer:} A
\\ \textit{Options:} A) All trapezoids are rectangles because they have at least one pair of parallel sides.; B) All pentagons are heptagons because they have more than 4 sides.; C) All hexagons are triangles because they have at least 3 sides.; D) All triangles are polygons because they have less than 4 sides.; E) All rectangles are squares because they have 4 right angles.
\\ \textit{Note:} The statement is incorrect because, while all rectangles are quadrilaterals with two pairs of parallel sides, not all trapezoids, which only require one pair of parallel sides, can be classified as rectangles.
\item \textbf{Answer:} C
\\ \textit{Options:} A) 10; B) 30; C) 23; D) 15; E) 18
\\ \textit{Note:} The correct answer is 23 because it is derived from applying Burnside's lemma to count the distinct colorings of the cube's vertices under the action of its rotational symmetries, which identifies the number of unique configurations that can be achieved with two colors.
\item \textbf{Answer:} C
\\ \textit{Options:} A) 2200; B) 1800; C) 2041; D) 2100; E) 1700
\\ \textit{Note:} The correct answer of 2041 meters is obtained by determining the maximum height of the rocket's trajectory, which occurs at the vertex of the parabolic function described by the vertical component of the trajectory c(t) = 200 t - 4.9 t², calculated using the formula for the vertex of a parabola, t = -b/(2 a), yielding the maximum height when substituting this t value back into the height equation.
\item \textbf{Answer:} D
\\ \textit{Options:} A) 1:03; B) 60:15; C) 4:01; D) 1:04; E) 3:01
\\ \textit{Note:} The correct answer is 1:4 because the ratio of the number of toucans (15) to the number of parrots (60) simplifies to 1:4 when both numbers are divided by their greatest common divisor, which is 15.
\end{enumerate}

\section*{True / False}
\begin{enumerate}[label=\arabic*.]
\setcounter{enumi}{5}
\item \textbf{Answer:} True
\\ \textit{Options:} A) True; B) False
\\ \textit{Note:} The magnitude of the vector v = (5, 0, 7) is calculated using the formula √(x² + y² + z²), which gives √(5² + 0² + 7²) = √(25 + 0 + 49) = √74, approximately equal to 8.6 units, making the statement false.
\item \textbf{Answer:} False
\\ \textit{Options:} A) True; B) False
\\ \textit{Note:} The statement is false because the upward velocity of the satellite at t = 1 is not 6.723 units per time interval, indicating that either the calculated value is incorrect or the conditions at that specific time do not yield that velocity.
\item \textbf{Answer:} False
\\ \textit{Options:} A) True; B) False
\\ \textit{Note:} The statement is false because the sum of two independent standard normal random variables, W(1) and W(2), follows a normal distribution with a mean of 0 and a variance of 2, so the probability that their sum exceeds 2 is not equal to 0.500, but rather significantly less than that.
\item \textbf{Answer:} False
\\ \textit{Options:} A) True; B) False
\\ \textit{Note:} The statement is false because two vectors in $R^2$ are linearly independent only if they are not scalar multiples of each other, which can occur if they are parallel or pointing in the same or opposite directions.
\item \textbf{Answer:} False
\\ \textit{Options:} A) True; B) False
\\ \textit{Note:} The statement is false because the integral $\int_{0}^\pi \left(\frac{\sin(123 x/2)}{\sin(x/2)}\right)^2 dx$ evaluates to a value that is not equal to 456.7890123, as the actual computation yields a different result based on the properties of the integral and the sine functions involved.
\end{enumerate}

\section*{Long Form Response}
\begin{enumerate}[label=\arabic*.]
\setcounter{enumi}{10}
\item \textbf{Answer:} \begin{align*} QED \&= (5+2 i)(i)(5-2 i)\\ \&=i(25-(2 $i)^2$)\\ \&=i(25+4)\\ \&=\ensuremath{\boxed{29 i}}. \end{align*}
\\ \textit{Note:} correct because it accurately applies the properties of complex numbers, specifically using the formula for the product of conjugates and the rules of multiplication involving the imaginary unit \( i \), leading to the correct result of \( 29 i \).
\item \textbf{Answer:} To express the quadratic function \(3 x^2 + x - 4\) in vertex form \(a(x - h)^2 + k\), we first complete the square. We factor out the coefficient of \(x^2\): \(3(x^2 + \frac{1}{3}x) - 4\). Next, we complete the square inside the parentheses: \(x^2 + \frac{1}{3}x = \left(x + \frac{1}{6}\right)^2 - \frac{1}{36}\). Thus, we rewrite the function as \(3\left(\left(x + \frac{1}{6}\right)^2 - \frac{1}{36}\right) - 4\), which simplifies to \(3\left(x + \frac{1}{6}\right)^2 - \frac{49}{12}\). In this vertex form, \(k = -\frac{49}{12}\) represents the minimum value of the function, indicating the lowest point on the graph, and it is significant as it determines the vertical position of the vertex, influencing the overall shape and direction of the parabola.
\\ \textit{Note:} correct because it demonstrates the process of completing the square to convert the standard form of the quadratic function into vertex form, accurately identifying the vertex coordinates and the significance of the constant \(k\) as the minimum value of the function, which corresponds to the vertical position of the vertex on the graph.
\end{enumerate}

\vfill
\noindent\textit{End of Answer Key}
\end{document}